% ============================================================================
% CHAPTER 14 TEACHER'S MANUAL: Building Fair and Ethical Systems
% Pedagogical guide for instructors
% ============================================================================

\section*{Chapter 14 Teacher's Manual: Building Fair and Ethical Systems}
\addcontentsline{toc}{section}{Chapter 14 Teacher's Manual}

% ============================================================================
\subsection*{Instructor Overview}
% ============================================================================

Chapter 14 is among the most sensitive chapters in the book. The subject matter---algorithmic bias, racial disparity in criminal justice algorithms, and the limits of human oversight as an ethics solution---can provoke strong emotional responses alongside intellectual ones. This is not a reason to avoid the material; it is a pedagogical opportunity. Students who feel personally invested in questions of fairness bring authentic engagement that moves discussion beyond the abstract.

\textbf{Key pedagogical goals:}
\begin{enumerate}
    \item Help students distinguish \emph{bias} (a value-neutral statistical term) from \emph{prejudice} (an intentional moral failure)---the chapter's most important conceptual move
    \item Develop comfort with mathematical impossibility results as legitimate intellectual conclusions, not defects in the argument
    \item Build the capacity to evaluate competing fairness claims without assuming one is simply ``right''
    \item Identify the structural conditions under which human oversight becomes meaningful versus performative
\end{enumerate}

\textbf{Prerequisites:} Students should have read Chapter 5 (annotation and labeling) and Chapter 7 (calibration and uncertainty). The technical appendix material on fairness metrics is optional for non-technical courses.

% ============================================================================
\subsection*{Learning Objectives}
% ============================================================================

By the end of this chapter, students should be able to:

\begin{enumerate}
    \item Identify and distinguish the five layers where bias enters AI systems (collection, annotation, model, deployment, feedback loops)
    \item State the fairness impossibility theorem and explain its normative implications
    \item Evaluate competing fairness metrics claims using the COMPAS case as a test
    \item Distinguish meaningful accountability structures (explainability, auditability, contestability) from accountability theater
    \item Articulate why well-intentioned human oversight does not automatically produce fair outcomes
    \item Apply a structured fairness audit framework to a real or hypothetical system
\end{enumerate}

% ============================================================================
\subsection*{Discussion Questions}
% ============================================================================

\subsubsection*{Opening Discussion (15 minutes, pre-lecture)}

\textbf{Q0:} Before reading, think of a time when a system or institution treated you, or someone you know, unfairly---either through a deliberate policy or through an impersonal process. What made it feel like a systemic failure rather than an individual mistake? What would have needed to change to fix it?

\textit{Instructor note:} This question frames the chapter's subject matter in terms of students' lived experience before introducing technical vocabulary. It often reveals that students' intuitions about fairness are already more sophisticated than they realize.

\subsubsection*{On the Layers of Bias}

\textbf{Q1:} The chapter identifies five layers where bias enters AI systems: collection, annotation, model, deployment, and feedback loops. Which is hardest to detect after the fact? Which is hardest to prevent? Are these the same layer?

\textbf{Q2:} ``Bias is not a bug you remove---it is a signal the model correctly learned from biased data.'' What are the implications of this statement for how organizations should approach bias remediation? What does it suggest about ``debiasing'' as a concept?

\textbf{Q3:} Amazon's recruiting AI adapted when engineers removed gender-correlated features---it found new proxy variables. What does this pattern of proxy adaptation suggest about the limits of feature engineering as a fairness strategy?

\subsubsection*{On Feedback Loops}

\textbf{Q4:} The predictive policing feedback loop shows arrests increasing in over-policed areas, which the model interprets as evidence that crime is higher there, generating more policing. The humans reviewing the system see more arrests and conclude the system is ``working.'' What cognitive bias among the human reviewers contributes to this outcome? What would a human reviewer need to know to recognize they are watching a feedback loop rather than model validation?

\textbf{Q5:} Feedback loops where model predictions influence future training data appear in many domains beyond policing: credit scoring, hiring, content recommendation, healthcare resource allocation. Choose a domain you know well and trace a plausible feedback loop. What would ``break the loop'' in your domain?

\subsubsection*{On the COMPAS Case and Fairness Impossibility}

\textbf{Q6:} ProPublica found that COMPAS had disparate false positive rates; Northpointe found that COMPAS was calibrated. Both were mathematically correct. Before reading the explanation: which claim seems more important to you, and why? After reading the mathematical proof of impossibility: has your intuition changed, or been reinforced?

\textbf{Q7:} The fairness impossibility theorem says the choice between fairness definitions is normative, not technical. Who should make this choice for a recidivism prediction tool used in sentencing decisions? Judges? Legislatures? Advocacy groups? The company that built the tool? What process would be appropriate?

\textbf{Q8:} The chapter argues that the COMPAS algorithm ``made the choice visible'' rather than ``making the choice.'' Is this defense satisfying? Does making an implicit social choice explicit (via an algorithm) change the moral responsibility of the people who deploy it?

\subsubsection*{On the Limits of Human Oversight}

\textbf{Q9:} The chapter identifies three reasons why adding human review doesn't automatically solve fairness problems: (1) humans absorb and legitimize AI bias, (2) humans carry their own biases, (3) humans can't track systemic effects while doing case review. Which of these three seems most practically tractable to address? What specific intervention would address each?

\textbf{Q10:} ``Human oversight, in this case, was not a check on the algorithm's failures. It was cover for them.'' What organizational conditions make this outcome---humans providing legitimizing cover rather than genuine oversight---most likely? What conditions would make genuine oversight more likely?

\subsubsection*{On Accountability Design}

\textbf{Q11:} The chapter distinguishes explainability (why did this prediction occur?), auditability (what patterns appear across many predictions?), and contestability (can individuals meaningfully challenge decisions?). Rate these in order of importance for each application: (a) medical AI making cancer screening recommendations, (b) a hiring algorithm screening résumés, (c) a recidivism prediction tool used in sentencing, (d) a social media content moderation algorithm.

\textbf{Q12:} The EU GDPR establishes a ``right to explanation'' for automated decisions. What would a genuinely informative explanation look like for each application in Q11? What would a formally compliant but practically useless explanation look like?

\subsubsection*{Synthesis and Application}

\textbf{Q13:} Design a fairness audit for a specific high-stakes AI system you care about (healthcare, hiring, lending, criminal justice, benefits eligibility). Specify: (a) which fairness definition(s) you would prioritize and why, (b) what data you would need, (c) who should conduct the audit, and (d) what you would do if the audit found disparate impact.

\textbf{Q14:} The chapter concludes that ``fairness is not a property a system achieves once---it is an ongoing relationship between a system, a deployment context, and an affected community.'' What institutional structures would be needed to operationalize this claim?

% ============================================================================
\subsection*{Classroom Activities}
% ============================================================================

\subsubsection*{Activity 1: The Hiring Algorithm Design Challenge (30 minutes)}

Working in groups of 4--5, students design a hiring algorithm for a software engineering role at a hypothetical company with a historically homogeneous workforce.

\textbf{Task sequence:}
\begin{enumerate}
    \item (5 min) List what data you would use to predict job performance. What proxies might encode demographic characteristics even if you don't include them explicitly?
    \item (8 min) Choose a fairness definition for your system. What does it mean concretely? Who are you protecting from what type of error?
    \item (8 min) A critic argues your chosen fairness definition ignores a competing valid concern. Respond. Can you satisfy both?
    \item (9 min) What human review process would you build into your system? Who reviews what? What authority do reviewers have?
\end{enumerate}

\textbf{Debrief questions:}
\begin{itemize}
    \item Which fairness definition did your group choose, and why?
    \item Where did your group most sharply disagree?
    \item What would your system need that you couldn't provide?
\end{itemize}

\subsubsection*{Activity 2: Role-Play Scenario---The Parole Board (25 minutes)}

\textbf{Roles assigned:}
\begin{description}
    \item[Role A] Defendant's attorney
    \item[Role B] Prosecutor
    \item[Role C] Parole board member
    \item[Role D] COMPAS system spokesperson
    \item[Role E] Researcher (can be played by multiple students)
\end{description}

\textbf{Scenario:} The parole board is reviewing the case of a defendant who received a high COMPAS score (classified high-risk for recidivism) but has served three years without incident in prison, has family support, and has completed vocational training.

Each role argues for their interpretation of the COMPAS score's relevance and appropriate weight in the decision. The parole board member decides.

\textbf{Debrief:} What information would have changed the decision? Was the COMPAS score more useful as a reason for the decision, or as a reason to scrutinize the decision? What does genuine contestability require in this scenario?

\subsubsection*{Activity 3: The Audit Simulation (45 minutes, computer lab)}

Using the Fairlearn library or a simplified spreadsheet version, students work with a synthetic dataset of hiring decisions.

\textbf{Tasks:}
\begin{enumerate}
    \item Compute overall accuracy, precision, and recall for a provided classifier
    \item Compute demographic parity difference and equalized odds difference by race and gender
    \item Identify which subgroups have the largest performance gaps
    \item Choose a mitigation strategy and evaluate its effect on both fairness metrics and overall accuracy
    \item Present findings: what tradeoff did your mitigation create?
\end{enumerate}

% ============================================================================
\subsection*{Assessment Rubrics}
% ============================================================================

\subsubsection*{Short Essay (500--800 words): ``Choose a Fairness Definition''}

\textit{Prompt: Select one high-stakes AI application. Argue for a specific fairness definition for that application. Acknowledge the strongest objection to your choice and respond to it.}

\begin{center}
\begin{tabular}{p{3cm}p{2.8cm}p{2.8cm}p{2.8cm}p{2cm}}
\toprule
\textbf{Criterion} & \textbf{Excellent (4)} & \textbf{Proficient (3)} & \textbf{Developing (2)} & \textbf{Beginning (1)} \\
\midrule
Fairness definition accuracy & Correctly states math meaning; connects to application & Mostly correct; minor imprecision & Some confusion between definitions & Conflates definitions \\
\addlinespace
Normative argument & Clear, grounded argument for why this definition matches application & Reasonable argument, some gaps & Partially developed & No clear argument \\
\addlinespace
Engagement with impossibility & Explicitly addresses tradeoff; explains what is sacrificed & Acknowledges tradeoff exists & Mentions but doesn't engage & Ignores impossibility \\
\addlinespace
Counterargument & Steelmans opposing view; responds substantively & Identifies opposing view; partial response & Weak or strawman & No counterargument \\
\bottomrule
\end{tabular}
\end{center}

% ============================================================================
\subsection*{Common Misconceptions to Address}
% ============================================================================

\textbf{Misconception 1:} ``Bias in AI is just prejudice by programmers.''

\textit{Correction:} Bias in the statistical sense means a systematic deviation from an unbiased standard. It enters through historical data that reflects past and present social inequalities---data that a model correctly learns from. The model is not prejudiced; the world it was trained on was.

\textbf{Misconception 2:} ``You can fix bias by removing sensitive attributes from the model.''

\textit{Correction:} The Amazon recruiting case illustrates the problem: models learn to use proxy variables. Removing race from a model trained on historical residential data does not prevent the model from learning neighborhood characteristics that correlate with race due to historical housing discrimination.

\textbf{Misconception 3:} ``Adding more human review solves the fairness problem.''

\textit{Correction:} The chapter's core argument is that this intuitive response is inadequate. Human reviewers carry their own biases, can be affected by automation bias, and cannot detect aggregate patterns while reviewing individual cases. Human oversight is necessary but not sufficient; it requires specific design to constitute real accountability.

\textbf{Misconception 4:} ``The fairness impossibility theorem means we can't have fair AI.''

\textit{Correction:} It means we can't simultaneously satisfy all possible definitions of fairness. This is a result about the structure of the problem, not a counsel of despair. It means we must choose which type of fairness matters most for a given application---and be honest about what we're giving up.

% ============================================================================
\subsection*{Connections to Other Chapters}
% ============================================================================

\begin{description}
    \item[Chapter 2] Calibration by demographic group extends Chapter 2's calibration concepts: a model can be well-calibrated overall but poorly calibrated for a minority subgroup.
    \item[Chapter 5] Annotation bias is a direct extension of the human judgment issues in Chapter 5; demographic homogeneity of annotator pools is a fairness issue as well as a quality issue.
    \item[Chapter 15] Feedback loop amplification from Chapter 14 is a variant of the systematic failure modes in Chapter 15; the audit trail requirements in Chapter 15 are prerequisites for the fairness monitoring discussed here.
    \item[Chapter 18] The civic participation theme in Chapter 18 connects directly to the question of who should make normative choices about fairness definitions.
\end{description}

% ============================================================================
\subsection*{Sample 50-Minute Lesson Plan}
% ============================================================================

\begin{center}
\begin{tabular}{p{2cm}p{5cm}p{4cm}}
\toprule
\textbf{Time} & \textbf{Activity} & \textbf{Materials} \\
\midrule
0:00--0:08 & Opening Q0 + Amazon story & Chapter reading \\
0:08--0:18 & Five layers of bias: lecture + Q1--Q3 & Slides \\
0:18--0:28 & COMPAS case: Q6 (before/after) & Fairness metrics handout \\
0:28--0:38 & Activity 2 (abbreviated): role-play parole board & Role cards \\
0:38--0:45 & Impossibility theorem + Q7 (who decides?) & Impossibility proof summary \\
0:45--0:50 & Misconceptions + exit ticket & --- \\
\bottomrule
\end{tabular}
\end{center}

\subsubsection*{75-Minute Extended Session}

\textbf{Add:} Activity 1 (Hiring Algorithm Design Challenge, 20 minutes condensed) between the five-layers section and the COMPAS case. This gives students a concrete design experience before encountering the impossibility result, which lands harder after they've tried and failed to satisfy multiple fairness criteria simultaneously.
